\SetPractica{Disolución de nitritos: Preparación}{Preparación de disolución patrón y estándar de trabajo de disolución de nitritos}

\section{Introducción teórica}



\section{Objetivos}

\begin{itemize}
    \item 
\end{itemize}

\section{Materiales, equipos y reactivos}

\begin{table}[H]
\centering{\small\setstretch{1.25}
\begin{tabular}{|m{0.325\linewidth}|m{0.325\linewidth}|m{0.325\linewidth}|} \hline
    
    \thead{Equipos} & \thead{Materiales} & \thead{Reactivos} \\ \hline
    
    {Desecador \par}
    {Balanza analítica \par}
    {Campana de extracción de gases \par}
    {Placa de calentamiento con agitador magnético \par}
    {Imán de agitación}
    
    &
    
    {Cápsula de porcelana \par}
    {Vasos de precipitado \par}
    {Matraz de aforación \par}
    {Probetas \par}
    {Embudos \par}
    {Pipetas graduadas}
    
    & 
    
    {Agua destilada \par}
    {Ácido clorhídrico (\ce{HCl}) \par}
    {Sulfanilamida \par}
    {Diclorhidrato de NEDA \par}
    {Acetato de sodio \par}
    {Nitrito de sodio \par}
    {Cloroformo}
    
    \\ \hline
\end{tabular}}\end{table}

\section{Procedimientos}

\subsection{Reactivo de color-buffer}
\begin{enumerate}
    \item A 250mL de agua destilada añadir
    \item 105mL de ácido clorhídrico concentrado
    \item 5g de sulfanilamida
    \item 0.5g de diclorhidrato de N-(1-naftil) etilendiamina.
    \item Agitar hasta disolver.
    \item Añadir 136g de acetato de sodio y agitar hasta disolver.
    \item Aforar a 500mL con agua destilada.
    \item Esta disolución es estable por varias semanas si se almacena en la oscuridad y en refrigeración. 
\end{enumerate}

\subsection{Disolución de referencia madre}
Concentración de masa de nitritos (\ce{NO2-}) = 100 mg/L de \ce{NO2-}
\begin{enumerate}
    \item Pesar con precisión 0.1493g de nitrito de sodio (\ce{NaNO2}), previamente desecado 24h en un desecador
    \item Disolver en agua y llevar a volumen de 1L.
    \item Preservar con 2mL de cloroformo por litro.
    \begin{enumerate}
        \item 1mL = 0.1mg de \ce{NO2-}.
    \end{enumerate}
\end{enumerate}

\subsection{Disolución de referencia de trabajo de nitritos}
Concentración de masa (\ce{NO2-}) = 1 mg/L de \ce{NO2-}
\begin{enumerate}
    \item Diluir 10mL de la disolución de referencia madre (6.8) a 1L.
    \item 1mL = 0.001mg de \ce{NO2-}
\end{enumerate}

\textbf{Notas:}
\begin{itemize}
    \item Cuando dice pesar aproximadamente y con precisión (patrones primarios) se refiere a pesar en balanza analítica.
    \item Cuando solo dice pesar aproximadamente se habla de pesar en balanza digital con 2 dígitos o uno si no se cuenta con dos dígitos después del punto.
    \item Cuando se habla de agua, siempre se refiere a agua desionizada.
    \item Cada disolución se conserva en frasco ámbar y en refrigeración, perfectamente etiquetada.
\end{itemize}

